\begin{textsample}{POS Dim 6 – global\_north\_2025\_04 – Score 4.00 – global\_north\_2025\_04/124017585.txt}  \label{ex:f6_pos_200}
All those who have stood by as Israel exterminates Gaza are guilty As Palestinians suffer, the world refuses to hold Israel accountable for violating international humanitarian law A Palestinian child eats his portion of a hot \textbf{meal} at a food distribution point in Gaza ' s Nuseirat camp on 19 April 2025 ( Eyad Baba/AFP ) Israel ' s genocide in Gaza, now into its 19th month, aims to destroy all life and dismantle the very \textbf{foundations} of survival in the besieged territory.
People in Gaza may die by missile strikes, disease, starvation or the collapsed healthcare system, as anguish weakens the body until illness overtakes it.
Every path leads to death in Gaza.
Death has become so frequent and familiar that Palestinians might simultaneously face multiple types of demise, as with the queues forming outside charity kitchens in an effort to stave off starvation, only to be bombed by Israeli aircraft while they wait.
This happened most recently in the southern al-Mawasi area, where one volunteer recalled to Middle East Eye @ @ @ @ @ @ @ @ @ @... and in the blink of an eye, everything turned upside down.
Blood filled the place, screams and cries of children and women rose, the food-filled pots were scattered, and the kitchen turned into a ball of fire."Since the start of the war in October 2023, Israel has targeted 26 food kitchens and 37 aid distribution centres in Gaza, according to local media reports.
New MEE newsletter : Jerusalem Dispatch Sign up to get the latest insights and analysis on Israel-Palestine, alongside Turkey Unpacked and other MEE newsletters While these kitchens often provide poor-quality food, typically limited to staples like lentils, beans and \textbf{rice}, Palestinians are willing to wait for hours in long queues just to get a plateful - a testament to the depth of Gaza ' s humanitarian catastrophe.
Many families have lost the ability to secure even a bare minimum of food.
These crowded kitchens have become their last refuge to ward off hunger.
Crippling \textbf{blockade} The crowding itself has produced new tragedies @ @ @ @ @ @ @ @ @ @ a boiling pot of food at a shelter in Nuseirat camp in January.
He suffered serious burns and later died of his injuries.
Food kitchens have spread across Gaza since the start of Israel ' s war of extermination.
They are run by volunteers and charitable organisations aiming to alleviate widespread poverty and hunger, but their ability to provide this service has been stymied by Israel ' s crippling \textbf{blockade} on Gaza, including restrictions on the entry of basic food supplies.
Since the start of March, Israel has not allowed a single food truck to enter Gaza.
All bakeries have been forced to close amid a lack of fuel and flour.
Most aid groups have suspended or cut their services.
The territory ' s two million people are facing a catastrophe.
At the same time, tens of thousands of children in Gaza are severely malnourished, threatening their future development."Children are eating less than one \textbf{meal} a day and struggling to get to the next,"Bushra al-Khalidi, Oxfam ' s regional policy @ @ @ @ @ @ @ @ @ @ canned food only...
Malnutrition and famine-like conditions are undoubtedly spreading in Gaza."There is no justification for using starvation as a tool of political pressure against an entire population - men, women and children alike.
Yet Israel pursues this policy with impunity, as far-right ministers intensify their genocidal rhetoric about not allowing"a gram of food or aid"into Gaza.
Israel has seen that the world lacks any serious will to hold it accountable for violating international law and humanitarian principles.
Had its allies taken a firm stand against its policy of collective punishment in Gaza, Israel would have been forced to retreat to some extent, fearing harm to its interests and international relationships.
But this never happened.
Israel is therefore not the only party responsible for the starvation of Gaza.
The culpability is shared among all those who continue to support Israel in any form, whether economically, politically or militarily.
And it is shared by all those who remain silent, giving a green light to the @ @ @ @ @ @ @ @ @ @.
The views expressed in this article belong to the author and do not necessarily reflect the editorial policy of Middle East Eye.
Ahmed Abu Artema is a Palestinian journalist and peace activist.
Born in Rafah, in 1984, Abu Artema is a refugee from Al Ramla village.
He authored the book"Organized Chaos".
Middle East Eye delivers independent and unrivalled coverage and analysis of the Middle East, North Africa and beyond.
To learn more about republishing this content and the associated fees, please fill out this form.
More about MEE can be found here.

% matched lemmas: blockade, foundation, meal, rice
\end{textsample}
